\subsubsection{Énergie des tuiles et réemploi des indices}\label{sec:energie-accelerateur}
Sur le chromosome~21, une tuile de profondeur \(h\) remplace \(h\) lancements par un seul, au prix du recalcul de son bord gauche. La figure~\ref{fig:energieaccelerateur} mesure cet échange sur deux régions homologues de GRCh38 et CHM13. Pour \(K\) appels de durée totale \(T\), le banc mesure l'énergie \(E\) du GPU par différence entre deux lectures du compteur de la bibliothèque NVML de NVIDIA~\cite{nvmldocs}.\footnote{La fonction \texttt{nvmlDeviceGetTotalEnergyConsumption} renvoie l'énergie cumulée du GPU en millijoules depuis le chargement du pilote. Les cœurs Grace et le reste de l’installation sont hors du périmètre de ce compteur.} La puissance moyenne \(P\) et le rendement \(\eta\) s'écrivent
\[
 P=E/T,\qquad \eta=\frac{Kmn}{10^9 TP}=\frac{Kmn}{10^9 E}\quad[\mathrm{GCUPS/W}].
\]
Les tuiles recalculent certaines cellules du halo. Le numérateur compte chaque cellule du problème une seule fois. Les cinq blocs suivent des ordres aléatoires. Chaque processus initialise CUDA et effectue un appel d'échauffement avant la mesure.

\chart{gpu/energie_gh200.png}{Le panneau a) donne le rendement moyen, puis l'écart-type et la demi-largeur de l'IC à 95\,\%. Le panneau b) compare les indices directs et compactés à profondeur seize. Les traits épais portent l'écart-type des cinq rapports appariés, les traits fins leur IC à 95\,\%.}{fig:energieaccelerateur}{python/studies/energy_gpu.py}

Sur la fenêtre colinéaire de 65\,536 bases, Nsight Systems compte 4\,097 lancements à seize lignes, initialisation comprise, contre 2\,049 à trente-deux. Le temps cumulé des noyaux passe pourtant de 49,3 à 51,8\,ms. La copie étalon atteint \(3{,}488\pm0{,}008\)\,To/s en comptant les octets lus et écrits (moyenne et écart-type). Le pilote réserve les compteurs HBM aux comptes autorisés. Pour des lots de deux puis huit secondes, l'énergie par appel vaut \(12{,}57\pm0{,}21\) et \(12{,}78\pm0{,}18\)\,J sur cette fenêtre.

Pour chaque symbole \(s\), la variante compacte stocke \(j-k_s(j)\) sur huit bits. Le code 255 renvoie à l'indice entier lorsque le symbole manque ou que l'écart atteint 255. Les quatre octets d'une colonne entrent dans un mot transféré en mémoire partagée au début de la tuile. L'insertion longue compte 2\,379 codes de repli sur 2\,149\,236 entrées. Son score vaut 13\,706 avec les deux noyaux. Le rendement dépend donc à la fois de la profondeur et de l'organisation des lectures.

\synthese{Sur l'insertion de 524\,288 bases, le compactage ramène le temps moyen de 1,878 à 1,461\,s et l'énergie du GPU de 534,0 à 427,7\,J par alignement. Le rapport énergétique apparié vaut 1,249, avec un IC à 95\,\% de 1,228 à 1,269.}
